\documentclass[11pt]{article}
\usepackage[margin=1in]{geometry}
\usepackage{amsmath,amssymb,amsthm}
\usepackage{booktabs}
\usepackage[hidelinks]{hyperref}

\newtheorem{theorem}{Theorem}
\newtheorem{lemma}[theorem]{Lemma}
\newtheorem{corollary}[theorem]{Corollary}
\theoremstyle{definition}
\newtheorem{definition}[theorem]{Definition}
\newtheorem{remark}[theorem]{Remark}

% The proved bound. Raise it only after docs/RESULTS.md records the extension as COMPUTATIONAL.
\newcommand{\NB}{100{,}000}
\newcommand{\NBexp}{100000}
\newcommand{\ncand}{1.0004\cdot 10^{15}}

\title{No multiplicative persistence $12$ below $10^{\NBexp}$:\\ a sieve over chain numbers\\[4pt]
\large Draft of 20 October 2026 --- computer-assisted, not peer reviewed}
\author{[Author]\thanks{Computations and drafting were done with the assistance of an AI system (Claude, Anthropic).
Code, logs and validation records: \url{https://github.com/CorvusCorpus/persistence-research}.}}
\date{}

\begin{document}
\maketitle

\begin{abstract}
Let $p(n)$ be the product of the decimal digits of $n$. The multiplicative persistence of $n$ is the number of
applications of $p$ needed to reach a single digit. The record value $11$ is attained by $277777788888899$, and it is
conjectured that no number has persistence $12$. The best published lower bound for the smallest number of
persistence $12$, if it exists, was $2.67\cdot 10^{30000}$ (T.~Peters, 2023). We show that no number below
$10^{\NBexp}$ has persistence $12$ or more, so the smallest one, if it exists, exceeds $2.67\cdot 10^{\NBexp}$. The
proof reduces the question to \emph{chain numbers}: $7$-smooth numbers with no digit $0$ that do not contain both a $5$
and an even digit. Measured by a ``digit length'' adapted to the digit product, every chain number up to length
$\NB$ is one of $1{,}117$ known numbers, all below $10^{49}$. The search covers $\ncand$ candidates. It rests on four
elementary lemmas: a suffix collapse that turns the last $13$ digits of every candidate into one bit of a precomputed
bitmap, an exact $4^k$ count of admissible endings, an affine formula for the next $13$ digits, and the uniqueness of
an odd-digit $2$-adic ``thread'' for multiples of high powers of $5$. We describe the implementation and its
validation: exact-arithmetic comparisons, an independent reference implementation, rediscovery of known zero-free
numbers, and two complete re-executions of the original run up to length $40{,}000$ on different hardware with an
independent (GPU) implementation of the core stages, which then extended the search to length $\NB$.
A random-digit heuristic, calibrated against the complete data, is reported separately; it is not part of the proof.
\end{abstract}

\section{Introduction}

For a positive integer $n$ let $p(n)$ be the product of its decimal digits and let $\operatorname{per}(n)$ be the least
$k\ge 0$ with $p^{k}(n)<10$. The smallest numbers of persistence $0,1,\dots,11$ form OEIS A003001~\cite{oeis}, ending with
\[
\begin{aligned}
  277777788888899 &\to 4996238671872 \to 438939648 \to 4478976 \to 338688 \to 27648\\
                  &\to 2688 \to 768 \to 336 \to 54 \to 20 \to 0 .
\end{aligned}
\]
Sloane~\cite{sloane} asked whether persistence is bounded; the standing conjecture is that $11$ is the maximum in base
$10$. Orbits that end in an odd digit are settled by Brier, Clavier, Gutsche and Naccache~\cite{bcgn}; a $2$-adic bound
of Fonga~\cite{fonga} makes the nonzero even final digits a finite (if large) computation. The hard case is the final
digit $0$, which is where the record orbit ends: it amounts to understanding the $7$-smooth numbers without a digit
$0$, a problem of the same family as Erd\H{o}s's question whether $2^{86}$ is the last power of $2$ without a digit $0$.

Lower bounds for the smallest number $a(12)$ of persistence $12$ come from computation. Chaffin~\cite{chaffin}
found that every $7$-smooth number in $(10^{140},10^{20000})$ contains a digit $0$, which gives $a(12)>10^{20585}$
(Radcliffe~\cite{radcliffe}). Peters~\cite{peters} enumerated the smallest representative of every digit-product
class up to $30{,}000$ digits (about $3.6\cdot 10^{13}$ candidates) and checked that each digit product has a $0$ among
its last $306$ digits, giving $a(12)>2.67\cdot 10^{30000}$. We prove:

\begin{theorem}\label{thm:main}
No positive integer below $10^{\NBexp}$ has multiplicative persistence $12$ or more. Consequently, if $a(12)$ exists,
then $a(12)>2.67\cdot 10^{\NBexp}$.
\end{theorem}

The proof is computer-assisted (Section~\ref{sec:results}). Its new ingredient is a sieve over \emph{chain numbers}
(Section~\ref{sec:chain}) in which most candidates are rejected by a single bit lookup that tests their last $13$
digits (Section~\ref{sec:suffix}). Section~\ref{sec:validation} describes how the computation was checked,
Section~\ref{sec:bridges} what a counterexample above the bound would have to look like, and
Section~\ref{sec:heuristic} a heuristic that is kept apart from the proof.

\section{Chain numbers}\label{sec:chain}

\begin{definition}
A number is \emph{$7$-smooth} if all its prime factors are at most $7$. For $x=2^a3^b5^d7^c$ put
\[ m(x)=\tfrac{a}{3}+\tfrac{b}{2}+c+d, \]
the \emph{digit length} of $x$. A \emph{chain number} is a $7$-smooth $x\ge 10$ with no digit $0$ that does not contain
both a digit $5$ and an even digit. We write $G$ for the set of chain numbers.
\end{definition}

A chain number is never divisible by $10$, so it has one of the forms $2^a3^b7^c$ (form~(i)) or $3^a5^b7^c$ with $b\ge1$
(form~(iii)); in form~(iii) the last digit is $5$, so every digit is odd. (Splitting the search for zero-free digit
products into these two families is standard; see~\cite{mcelwee}.)

\begin{lemma}[chain lemma]\label{lem:chain}
Let $n$ have persistence $k\ge 12$ and put $c_i=p^{\,i+1}(n)$. Then $c_0,\dots,c_8$ are chain numbers, and
$m(c_0)\le \operatorname{len}(n)$, the number of digits of $n$.
\end{lemma}

\begin{proof}
By definition $c_i\ge 10$ for $i\le k-2$, and $c_{k-1}$ is a single digit. Every $c_i$ is a digit product, hence
$7$-smooth, and $c_i\ge10$ is nonzero. For $i\le k-3$, $c_{i+1}=p(c_i)\neq 0$, so $c_i$ has no digit $0$. If $c_i$ contained a
$5$ and an even digit, then $10\mid p(c_i)=c_{i+1}$, so $c_{i+1}$ would end in $0$ and $c_{i+2}=0$, which is impossible for
$i+2\le k-2$. Hence $c_i\in G$ for $i\le k-4$, in particular for $i\le 8$. For the second claim, a digit $d$ of $n$
multiplies $p(n)$ by $d$, which raises $m$ by $0,\tfrac13,\tfrac12,\tfrac23,1,\tfrac56,1,1,1$ for $d=1,\dots,9$; each
contribution is at most $1$.
\end{proof}

\begin{corollary}\label{cor:reduction}
Suppose every chain number $x$ with $m(x)\le N$ is one of a known finite list, all of whose members have
persistence at most $10$. Then every $n$ with at most $N$ digits has persistence at most $11$.
\end{corollary}

\begin{proof}
If $\operatorname{per}(n)\ge 12$ and $\operatorname{len}(n)\le N$, then $c_0=p(n)$ is a chain number with $m(c_0)\le N$
(Lemma~\ref{lem:chain}), hence on the list, and $\operatorname{per}(n)=1+\operatorname{per}(c_0)\le 11$.
\end{proof}

The bound $a(12)>2.67\cdot10^{N}$ follows from Corollary~\ref{cor:reduction} by the candidate argument of
Peters~\cite{peters}. Suppose $a=a(12)$ exists. Permuting digits, deleting a digit $1$, or replacing some digits by
fewer digits with the same product ($22\to4$, $23\to6$, $24\to8$, $33\to9$, $266\to89$) gives a number with the same
digit product, hence the same persistence, and a smaller one unless nothing changes; and a number containing a $5$
and an even digit has persistence at most $2$. So the digits of $a$ are nondecreasing and contain none of these
configurations. By Corollary~\ref{cor:reduction}, $a$ has at least $N+1$ digits. If its first digit is $2$, its second
is at least $6$, and after $2,6$ every further digit is at least $7$. Hence
$a\ge 2\,6\,7\cdots7$ (with $N+1$ digits) $>2.67\cdot10^{N}$.

\begin{remark}
A chain number with a $5$ has only odd digits. Per digit, a random number satisfies the chain condition with probability
about $0.8$, against $0.9$ for being zero-free; this, and the weight $m$, which charges each prime by the shortest digit
string that produces it, are what make chain numbers much rarer than zero-free $7$-smooth numbers: $1{,}117$ chain
numbers are known, the largest with $49$ digits, against $12{,}606$ zero-free ones up to $140$ digits.
\end{remark}

\section{Suffix structure}\label{sec:suffix}

Throughout, $K\ge1$, $P_K=4\cdot5^{K-1}$ is the order of $2$ modulo $5^K$, and $\mu_3,\mu_7$ are the discrete logarithms
$2^{\mu_3}\equiv3$, $2^{\mu_7}\equiv7\pmod{5^K}$ (they exist because $2$ is a primitive root modulo every power of $5$).

\begin{lemma}[suffix collapse]\label{lem:L1}
If $a\ge K$ then $2^a3^b7^c\equiv 2^e\pmod{10^K}$ for every $e\ge K$ with $e\equiv a+b\mu_3+c\mu_7\pmod{P_K}$.
\end{lemma}

\begin{proof}
Both sides are $\equiv0\pmod{2^K}$, and both are $\equiv 2^{a+b\mu_3+c\mu_7}\pmod{5^K}$; apply the Chinese remainder
theorem.
\end{proof}

So the last $K$ digits of every form-(i) number with $a\ge K$ are those of a power of two, determined by one residue
modulo $P_K$. For $K=13$, a bitmap of $P_{13}=976{,}562{,}500$ bits records which residues have last $13$ digits
free of $0$ and $5$. Along a \emph{row} ($b,c$ fixed, $a$ increasing) the residue increases by one at each step, so a
row is a contiguous range of bits.

\begin{lemma}[exact branching]\label{lem:L2}
For $k\ge1$, among the $P_k$ residues $e\bmod P_k$ ($e\ge k$), exactly $4^k$ give a power $2^e$ whose last $k$ digits avoid
$0$ and $5$.
\end{lemma}

\begin{proof}
For $k=1$ the last digit of $2^e$ runs through $2,4,8,6$. Fix $e\ge k+1$ and consider the five lifts $e_t=e+tP_k$,
$t=0,\dots,4$, of $e\bmod P_k$ to residues modulo $P_{k+1}$. Write $2^{e_t}\equiv x_0+10^k d_t\pmod{10^{k+1}}$ with
digits $d_t$; the last $k$ digits do not depend on $t$. Since $2^{P_k}=1+5^ku$ with $5\nmid u$, we have
$2^{e_t}\equiv 2^e(1+t\,5^ku)\pmod{5^{k+1}}$, so $2^k(d_t-d_0)\equiv 2^et u\pmod 5$ and the $d_t$ are distinct modulo
$5$. Modulo $2^{k+1}$ both $2^{e_t}$ and $2^{e_0}$ vanish, so $2^k5^k(d_t-d_0)\equiv0$ and the $d_t$ have equal parity.
Hence the $d_t$ are the five digits of one parity, and exactly one of them lies in $\{0,5\}$.
\end{proof}

The bitmap therefore has exactly $4^{13}=67{,}108{,}864$ set bits (a density of $0.8^{12}\approx0.0687$), and stage~1 of
the sieve passes exactly that fraction of a complete period. Lifting arguments of this kind are standard for digit
problems on powers of two (see, e.g.,~\cite{dumitru}); the exact count for the excluded set $\{0,5\}$ is what makes the
chain condition well suited to a bitmap.

\begin{lemma}[affine lift]\label{lem:L3}
Let $Q=5^K$, $P=P_K$, $0\le e<P$, and let $y=2^{e-2K}\bmod 5^{2K}=A_{\mathrm{lo}}+QA_{\mathrm{hi}}$ with
$0\le A_{\mathrm{lo}},A_{\mathrm{hi}}<Q$. Let $w=(2^{P}\bmod 5^{2K}-1)/Q$. Then for every $j\ge0$ with $e+jP\ge 2K$, the
digits $K+1,\dots,2K$ of $2^{e+jP}$ form the integer
\[ \Bigl\lfloor\frac{2^K A_{\mathrm{lo}}}{Q}\Bigr\rfloor+2^K\bigl((A_{\mathrm{hi}}+A_{\mathrm{lo}}\,w\,j)\bmod Q\bigr). \]
\end{lemma}

\begin{proof}
$2^{e+jP}\bmod 10^{2K}=2^{2K}\bigl(2^{e+jP-2K}\bmod 5^{2K}\bigr)$, and $2^{e+jP-2K}\equiv y(1+Qw)^j\equiv y(1+Qwj)\equiv
A_{\mathrm{lo}}+Q(A_{\mathrm{hi}}+A_{\mathrm{lo}}wj)\pmod{Q^2}$. Dividing $2^{2K}(A_{\mathrm{lo}}+QY)$, $Y<Q$, by $10^K=2^KQ$
gives the stated quotient.
\end{proof}

By Lemma~\ref{lem:L1} with $K$ and $2K$, the full exponent residue of a form-(i) number modulo $P_{2K}$ is $e+jP_K$ with
$j$ constant along a row. So after the bitmap, the next $K$ digits of each surviving point cost one modular
multiplication by the row constant $j$.

\begin{lemma}[odd thread]\label{lem:L4}
There is exactly one residue $r_{40}$ modulo $10^{40}$ with $5^{40}\mid r_{40}$ and all $40$ digits odd:
\[ r_{40}=\cdots 9313975971735999337397515773773193359375 . \]
Write $r_{40}\equiv -5^{t}\pmod{2^{40}}$ and $3\equiv-5^{f_3}$, $7\equiv-5^{f_7}\pmod{2^{40}}$. If $b\ge40$, the last $40$
digits of $x=3^a5^b7^c$ are all odd if and only if $a+c$ is odd and $b\equiv t-af_3-cf_7\pmod{2^{38}}$.
\end{lemma}

\begin{proof}
If $r$ has $k$ odd digits and $5^k\mid r$, the lifts $r+j10^k$ divisible by $5^{k+1}$ are $j\equiv j_0\pmod5$, i.e.
$j_0$ and $j_0+5$, one odd and one even; so exactly one lift survives, starting from $r=5$. For $b\ge40$,
$x\equiv0\pmod{5^{40}}$ and $x\equiv(-1)^{a+c}5^{\,b+af_3+cf_7}\pmod{2^{40}}$; since $5$ has order $2^{38}$ modulo $2^{40}$
and $-1$ is not a power of $5$ there, $x\equiv r_{40}\pmod{10^{40}}$ exactly under the stated conditions.
\end{proof}

So for each pair $(a,c)$ at most one residue class of $b$ modulo $2^{38}$ survives the last $40$ digits: the odd
form costs $O(1)$ per pair $(a,c)$ for $b\ge40$.

\section{The sieve}\label{sec:sieve}

\paragraph{Regions.} By Corollary~\ref{cor:reduction} it suffices to examine every $x=2^a3^b7^c$ with $2a+3b+6c\le6N$
and every $x=3^a5^b7^c$, $b\ge1$, with $a+2b+2c\le 2N$. All such $x<10^{131}$ (there are $4{,}483{,}467$) were enumerated
directly; among them are exactly the $1{,}117$ chain numbers of the known list~\cite{repo}. The sieve handles
$x\ge10^{130}$, so the two ranges overlap. Shells $N_0<m\le N$ are run separately with an extra lower bound.

\paragraph{Main path} (form (i), $a\ge26$, $x\ge10^{130}$). Each candidate passes through
\begin{enumerate}
\item the bitmap of Lemma~\ref{lem:L1}, $K=13$ (last $13$ digits): survival $0.8^{12}$;
\item the affine lift of Lemma~\ref{lem:L3} (digits $14$--$26$): survival $0.8^{13}$ of the rest;
\item a leading-digit test: $t=aL_2+bL_3+cL_7\bmod 2^{64}$ with $L_p=\operatorname{round}(2^{64}\log_{10}p)$ approximates
  $2^{64}\{\log_{10}x\}$ within $(a+b+c)/2$ units, and a $2^{18}$-bin table rejects only bins in which every point (with
  a margin of $2^{-40}$) has a $0$ or $5$ among its first $5$ significant digits;
\item exact digits $27$--$126$ by a base-$10^9$ limb product;
\item GMP: the last $600$ digits, then the whole number.
\end{enumerate}
Only stage 3 uses approximations, and it rejects a candidate only when the rejection is certain under a proved error
bound; everything else is exact integer arithmetic.

\paragraph{Direct and odd paths.} Form (i) with $a<26$ is stepped along $b$ with three $18$-digit limbs. Form (iii) with
$b\le39$ is stepped along $a$ with an $18$-digit filter, and $b\ge40$ uses Lemma~\ref{lem:L4}. Survivors are decided
with GMP.

\paragraph{Cost.} The region has about $N^3$ points, so the cost grows like $N^3$. Stages $1$--$2$ cost about
$0.17$\,ns per candidate on one modern CPU core, or about $5$\,ps on an RTX~3060~Ti.

\section{Implementation and validation}\label{sec:validation}

The program (\texttt{gsieve.c}, C with GMP, about 1,000 lines) has several interchangeable implementations of stages
1--2: AVX-512 with IFMA (the original production build), AVX-512, AVX2, a scalar reference, and CUDA.
\begin{itemize}
\item \emph{Arithmetic.} Every vector implementation divides by $10^4$ as $\lfloor vM/2^{45}\rfloor$ with
  $M=\lceil2^{45}/10^4\rceil$. This is exact for $0\le v<2^{45}/1168$ (Granlund--Montgomery~\cite{gm}), and the operands
  are split so that every divided value is below $2^{31}$; the identity was checked for all $v<2^{31}$. The lift uses
  Montgomery multiplication~\cite{montgomery} modulo $5^{13}$.
\item \emph{Mutual agreement.} All implementations make identical decisions. Dumps of every stage-2 survivor and
  stage-3/4 decision agree across implementations and thread counts, and the CUDA dumps are byte-identical to the CPU
  dumps.
\end{itemize}

The computation was checked as follows.
\begin{enumerate}
\item \emph{Exact arithmetic on full regions.} Every stage decision on full regions up to $N=150$ (both $K=12$ and
  $K=13$ builds) and on complete production slices (single values of $c$ at $N=36{,}000$ and in the shell, up to
  $3.6\cdot10^8$ candidates each) was compared with big-integer arithmetic.
\item \emph{Positive control.} In a zero-free mode the sieve must rediscover known numbers. End to end at $N=160$ it
  reproduces exactly the $314$ zero-free $7$-smooth numbers of that size from an independent list complete below
  $10^{5000}$.
\item \emph{Brute force.} The direct and odd paths agree with brute force on full small regions. The odd thread was also
  checked digit by digit on production regions that contain real thread candidates.
\item \emph{Independent reference.} A separate plain-Python implementation shares no code with the program and decides
  everything from the digits of $x$; it re-derives the constants, Lemma~\ref{lem:L2} for $k\le9$ and the odd thread. It
  agrees with the program on full regions and on production slices, including all odd-thread candidates found there.
\item \emph{Kernel tests (GPU).} On the whole table, $6.9\cdot10^{10}$ stage-2 decisions of the GPU agree with an
  independent host formula. On $2^{20}$ random points, the decisions agree with $2^E\bmod10^{26}$ computed by GMP.
\item \emph{Complete re-executions.} The logged production run for $m\le40{,}000$ ($6.4067\cdot10^{13}$ candidates, computed on a cloud machine with
  ECC memory) was re-executed twice on a desktop PC with the GPU implementation of stages 1--2. Every counter of every
  invocation and every per-$c$ count agreed with the original logs.
\item \emph{Statistics.} In every run the survival rates of stages 1 and 2 match $0.8^{12}=0.0687195$ and
  $0.8^{13}=0.0549756$. Stage 1 agrees to all printed digits; stage 2 is within $3$ binomial standard deviations
  (Table~\ref{tab:runs}).
\item \emph{Fault tolerance.} The shells beyond $m=40{,}000$ ran once, on a desktop without ECC memory. The GPU tables
  were re-verified every minute against the host copies, and output was written only after a passed check, so a log
  never contains results computed with corrupted tables. An independent second pass of these shells on the CPU,
  compared counter by counter, is planned; it would also cover transient faults.
\end{enumerate}

\section{Results}\label{sec:results}

\begin{table}[ht]
\centering\small
\begin{tabular}{lrrrr}
\toprule
run & main path & direct & odd & total \\
\midrule
$m\le 36{,}000$ & $46{,}626{,}195{,}781{,}376$ & $33{,}689{,}253{,}108$ & $50{,}490{,}005{,}607$ & $46{,}710{,}375{,}040{,}091$ \\
$36{,}000<m\le40{,}000$ & $17{,}337{,}008{,}704{,}000$ & $7{,}903{,}308{,}000$ & $11{,}850{,}072{,}000$ & $17{,}356{,}762{,}084{,}000$ \\
$40{,}000<m\le50{,}000$ & $60{,}979{,}301{,}760{,}000$ & $23{,}398{,}270{,}000$ & $35{,}085{,}180{,}000$ & $61{,}037{,}785{,}210{,}000$ \\
$50{,}000<m\le60{,}000$ & $90{,}974{,}701{,}760{,}000$ & $28{,}598{,}270{,}000$ & $42{,}885{,}180{,}000$ & $91{,}046{,}185{,}210{,}000$ \\
$60{,}000<m\le70{,}000$ & $126{,}970{,}101{,}760{,}000$ & $33{,}798{,}270{,}000$ & $50{,}685{,}180{,}000$ & $127{,}054{,}585{,}210{,}000$ \\
$70{,}000<m\le80{,}000$ & $168{,}965{,}501{,}760{,}000$ & $38{,}998{,}270{,}000$ & $58{,}485{,}180{,}000$ & $169{,}062{,}985{,}210{,}000$ \\
$80{,}000<m\le90{,}000$ & $216{,}960{,}901{,}760{,}000$ & $44{,}198{,}270{,}000$ & $66{,}285{,}180{,}000$ & $217{,}071{,}385{,}210{,}000$ \\
$90{,}000<m\le100{,}000$ & $270{,}956{,}301{,}760{,}000$ & $49{,}398{,}270{,}000$ & $74{,}085{,}180{,}000$ & $271{,}079{,}785{,}210{,}000$ \\
\midrule
$m\le 100{,}000$ & & & & $1{,}000{,}419{,}848{,}384{,}091$ \\
\bottomrule
\end{tabular}
\caption{Candidates examined (plus $4{,}483{,}467$ numbers below $10^{131}$ enumerated directly). No run produced a chain
number. The first two runs used AVX-512 CPUs (ECC memory) and were later re-executed on a GPU; the six shells beyond
$m=40{,}000$ ran on a GPU in $86$ minutes.}\label{tab:runs}
\end{table}

\begin{theorem}[computational]\label{thm:comp}
Every chain number $x$ with $m(x)\le\NB$ is one of the $1{,}117$ chain numbers below $10^{49}$; the largest is
$4613824623662469877463718941499166849983276122112=2^{132}\cdot3\cdot7^{10}$. All of them have persistence at most
$10$, and only $4996238671872=2^{19}3^47^6$ and $937638166841712=2^43^{20}7^5$ have persistence exactly $10$.
\end{theorem}

Theorem~\ref{thm:main} follows from Theorem~\ref{thm:comp} and Corollary~\ref{cor:reduction}.

\section{Above the bound: bridges}\label{sec:bridges}

\begin{lemma}\label{lem:bridge}
Assume Theorem~\ref{thm:comp} for $N=\NB$, and let $n$ have persistence at least $12$, with $c_i$ as in
Lemma~\ref{lem:chain}. Then either $c_0,\dots,c_8$ all have $m>\NB$, or some $c_t$ ($t\le7$) is a \emph{bridge}: a chain
number with more than $69{,}897$ digits whose digit product is one of the $1{,}117$ known chain numbers. A bridge has
at most $161$ digits different from $1$, so more than $99.77\%$ of its digits are $1$s.
\end{lemma}

\begin{proof}
By Theorem~\ref{thm:comp}, $m(c_0)>\NB$; let $t$ be the largest index $\le8$ with $m(c_t)>\NB$. If $t<8$ then
$c_{t+1}=p(c_t)$ is a chain number with $m\le\NB$, hence known and below $2^{162}$, and each digit of $c_t$ other than
$1$ contributes a factor at least $2$. Finally $m(x)\le\log_{10}x/\log_{10}5<1.431\log_{10}x$ for $7$-smooth $x$, so
$m(c_t)>\NB$ forces more than $69{,}897$ digits.
\end{proof}

A bridge is a near-repunit: writing it as $(10^L-1)/9+\sum_i (d_i-1)10^{a_i}$ with at most $161$ digits $d_i\neq1$
gives $S$-unit equations with up to $161$ terms, far beyond the reach of the linear-forms methods
of~\cite{bcgn,fonga}.

\section{A heuristic (not part of the proof)}\label{sec:heuristic}

Model the digits of a $7$-smooth number with $L$ digits as independent and uniform. Against the complete list below
$10^{5000}$ the model predicts $5{,}622$ zero-free numbers with $20$--$39$ digits (actual $5{,}581$), $2{,}059$ with $40$--$59$
digits (actual $2{,}111$), and $7.3$ chain numbers with $40$--$59$ digits (actual $7$). A number of persistence $12$
exists if and only if some $7$-smooth number has digit product $4996238671872$ or $937638166841712$ (unless an unknown
persistence-$10$ digit product exists beyond $10^{20000}$). Under the model, the expected number of such $7$-smooth
numbers is about $0.019$, concentrated at $18$--$22$ digits, a range long excluded. The expected number of chain numbers
beyond $25{,}000$ digits is below $10^{-2400}$. Larger computational bounds therefore certify an exponentially thin tail;
they are not evidence in either direction about the conjecture.

\section{Reproducibility}

The repository~\cite{repo} contains the source, the production logs ($\sim$13\,MB, plus compressed archives for
later shells), SHA-256 manifests, and all validation logs with the exact commands, build flags, binary hashes and
commits. The original run up to length $40{,}000$ took about $5.5$ core-hours on a $2.1$\,GHz AVX-512 machine; its
complete re-execution takes about $8$ minutes on an RTX~3060~Ti with a $20$-thread CPU, and the six shells up to length
$\NB$ ($9.36\cdot10^{14}$ candidates) took $86$ minutes there.

\begin{thebibliography}{99}
\bibitem{oeis} OEIS Foundation, Sequence A003001, \emph{The On-Line Encyclopedia of Integer Sequences},
  \url{https://oeis.org/A003001} (accessed 20 October 2026).
\bibitem{sloane} N.~J.~A.~Sloane, The persistence of a number, \emph{J. Recreational Math.} 6 (1973), 97--98.
\bibitem{bcgn} E.~Brier, C.~Clavier, L.~Gutsche, D.~Naccache, The multiplicative persistence conjecture is true for odd
  targets, arXiv:2110.04263 (2021).
\bibitem{fonga} P.~N.~Fonga, The multiplicative persistence conjecture: resolving the $2$-adic obstruction for nonzero
  even targets, arXiv:2608.27802 (2026).
\bibitem{chaffin} B.~Chaffin, comment in OEIS A003001, 29 September 2016.
\bibitem{radcliffe} D.~Radcliffe, comment in OEIS A003001, 22 March 2019.
\bibitem{peters} T.~Peters, comment in OEIS A003001, 19 September 2023.
\bibitem{mcelwee} K.~McElwee, An algorithm for multiplicative persistence research, blog post, 13 July 2019.
\bibitem{dumitru} B.~C.~Dumitru, Notes on the finiteness of powers of two with all even digits, arXiv:2503.23177 (2025).
\bibitem{gm} T.~Granlund, P.~L.~Montgomery, Division by invariant integers using multiplication, \emph{Proc. PLDI}
  (1994), 61--72.
\bibitem{montgomery} P.~L.~Montgomery, Modular multiplication without trial division, \emph{Math. Comp.} 44 (1985),
  519--521.
\bibitem{repo} Persistence research repository, \url{https://github.com/CorvusCorpus/persistence-research}
  (\texttt{docs/CHAIN\_SIEVE.md}, \texttt{docs/NOTEBOOK.md}).
\end{thebibliography}

\end{document}
